\newpage
\begin{abstract}
本文针对2020年CUMCM B题电动汽车充电站布局优化问题，建立了多目标空间优化模型。首先对城市人口分布、交通网络和现有充电设施数据进行分析，构建了充电需求强度指数；其次建立了考虑覆盖最大化与投资成本最小化的双目标选址-定容优化模型；最后设计了兼顾公平性的站点布局方案。

\textbf{针对问题一}，我们综合人口密度、路网密度、现有设施等因素构建了需求强度指数：
\I_i = w_1 \cdot \frac{Pop_i}{\max(Pop)} + w_2 \cdot \frac{Road_i}{\max(Road)} + w_3 \cdot \left(1 - \frac{Station_i}{\max(Station)}\right)
权重采用熵权法客观确定，避免了主观赋权的偏差，确保了评价结果的科学性和公正性。

\textbf{针对问题二}，我们建立了覆盖-成本双目标优化模型：
\\min \quad Z_1 = \sum_j C_j x_j \qquad \max \quad Z_2 = \sum_i D_i \cdot \mathbb{1}[\min_j d_{ij} x_j \leq r]
采用NSGA-II算法求解，种群大小80，迭代150代，得到28个非支配解，覆盖完整的Pareto前沿。

\textbf{针对问题三}，引入服务均衡度指标量化公平性：
\\text{Equity} = \sqrt{\frac{1}{n}\sum_{i=1}^{n}(C_i - \bar{C})^2}
在优化中将其纳入目标函数，确保城市中心和郊区的充电服务公平性，提升整体社会福祉。

本文方法可为城市充电基础设施规划提供科学的决策支持，具有广泛的应用前景。
\keywords{电动汽车} \textbf{充电站选址} \textbf{多目标优化} \textbf{NSGA-II} \textbf{覆盖模型}
\end{abstract}